% ===========================================================================
%  Capitolo 7 — Esperimenti
%
%  REGOLA: questo capitolo descrive il DISEGNO. I risultati stanno nel
%  capitolo 8. Se una frase riporta un valore misurato di una metrica, è nel
%  posto sbagliato.
% ===========================================================================
\chapter{Esperimenti}
\label{cap:esperimenti}

Questo capitolo descrive il disegno dei quattro esperimenti: per ciascuno,
l'ipotesi, il punto esatto di intervento, la baseline, il protocollo di
esecuzione e il conteggio dei parametri. Il criterio di rifiuto è comune a tutti
ed è quello enunciato nella \cref{sec:falsificazione}: un'ipotesi è
\rifiutata{} se l'intervallo di confidenza al \SI{95}{\percent} della differenza
appaiata contiene lo zero oppure giace interamente nella direzione contraria a
quella attesa. Le configurazioni complete sono nell'\cref{app:configurazioni}, i
comandi nell'\cref{app:riproduzione} e i risultati nel \cref{cap:risultati}.

\section{Esperimento 01 --- Proiezioni \texorpdfstring{$Q,K,V$}{Q,K,V} quantistiche}
\label{sec:esp-gpt}

\begin{ipotesi}[H1]
Sostituendo le tre proiezioni lineari $Q$, $K$, $V$ di ogni testa di attenzione
con circuiti quantistici variazionali, e mantenendo invariato ogni altro
componente, la modellazione del linguaggio migliora: $\PPL_q < \PPL_c$.
\end{ipotesi}

La metrica primaria è la validation loss, con la perplexity che ne deriva.

L'intervento è confinato in \code{Head.\_get\_layer}
(\cref{ssec:punto-sostituzione}): con \code{use\_quantum=False} il metodo
restituisce \code{nn.Linear(n\_embd, head\_size, bias=False)}, con
\code{use\_quantum=True} restituisce \classe{QuantumLayerAdapter}. La catena di
forme attraversata da una singola proiezione quantistica è
\begin{equation}
  \shape{\dbatch, \dseq, \demb}
  \xrightarrow{\text{adapter}} \shape{\dbatch, \dseq, \nq}
  \xrightarrow{\tanh\cdot\pi} [-\pi,\pi]^{\nq}
  \to \texttt{VQC}
  \to \shape{\dbatch, \dseq, \nq} \equiv \shape{\dbatch, \dseq, \dhead}.
  \label{eq:shapes-gpt-quantum}
\end{equation}

A ogni \emph{forward pass} si valutano $n_{\text{layer}} \times \nhead \times 3$
circuiti indipendenti. Il vincolo strutturale $\dhead = \nq$ della
\cref{eq:nqubits-gpt} determina quali configurazioni siano ammissibili, e il
metodo solleva \code{ValueError} quando non è soddisfatto: la coppia
$(\demb, \nhead)$ non può quindi essere scelta liberamente, ma solo in modo che il
quoziente sia un numero di qubit simulabile.

La baseline è lo stesso modello con \code{use\_quantum=False}. Restano costanti
corpus, tokenizer, preset di configurazione, seed, numero di iterazioni,
ottimizzatore, dimensione di batch e lunghezza di contesto.

Il repository espone sei preset, confrontati nella \cref{tab:preset-gpt}; in tutti
$\nqlayers = 2$ e \code{q\_device = "default.qubit"}. Il confronto principale usa
\code{thesis}, fissato dopo un pilot e prima della campagna definitiva: tokenizer
a caratteri, $\demb = 8$, $\nhead = 2$, $\nq = 4$, batch 8, contesto 16, 1500
iterazioni, valutazione ogni 100 iterazioni su 100 batch e dropout nullo. Il
preset \code{big} non è utilizzabile nella variante quantistica, perché 64 qubit
non sono simulabili.

\begin{limite}
Il preset \code{fast} ha \code{use\_quantum=True} come valore predefinito. In un
confronto appaiato il campo va sovrascritto in modo esplicito per la variante
classica. La modalità \code{compare} di \file{main.py} lo fa impostando il campo
per ciascuna delle due varianti prima di costruire il modello, ma la cosa va
segnalata perché un'esecuzione manuale delle due varianti con lo stesso comando
produrrebbe due modelli quantistici.
\end{limite}

Il confronto principale usa esclusivamente \file{data/input.txt}, il corpus
Shakespeare distribuito con il repository. Contiene \num{1115390} caratteri, e lo
split deterministico 90/10 ne assegna \num{1003851} al training e \num{111539}
alla validation; \classe{CharTokenizer} produce un vocabolario di 65 simboli.
L'hash SHA-256 del corpus è
\code{434c0554a8c4c53dc17e56a0abb0f30b88f83cbceb0289cb897db68c25e89eba}.

La modalità \code{compare} esegue entrambe le varianti per ciascun seed, con
l'appaiamento descritto nella \cref{sec:multiseed}, e scrive l'aggregato in
\file{experiments/comparison\_<nome>.json}. Il file di progresso è aggiornato dopo
ogni variante e \code{--resume} salta esclusivamente i run completi.

Quanto ai parametri, per il preset \code{thesis} la singola proiezione passa da 32
parametri classici a 60 quantistici (\cref{tab:parametri-proiezione}). Con
$2\times2\times3 = 12$ proiezioni la differenza sul modello completo è di
$12\times28=\num{336}$ parametri, cioè \num{2945} nella variante classica contro
\num{3281} in quella quantistica ($+11{,}4\,\%$).

Vengono registrati validation loss e perplexity per seed, conteggio dei parametri,
tempi di addestramento e per iterazione, configurazione completa e manifest di
provenienza. La generazione di testo è prodotta come materiale qualitativo:
nessuna conclusione quantitativa è tratta dai campioni generati.

\section{Esperimento 02 --- \emph{Patch embedding} quantistico nel ViT}
\label{sec:esp-vit}

\begin{ipotesi}[H2]
Sostituendo la proiezione lineare delle patch con un circuito quantistico
variazionale, e mantenendo invariato l'encoder, la classificazione di immagini
biomediche migliora: $\accte^{\,q} > \accte^{\,c}$.
\end{ipotesi}

Metrica primaria è la test accuracy; macro-AUC e macro-F1 sono secondarie.

In \classe{HybridQCNNViT} il campo \code{use\_quantum} seleziona il \emph{patch
embedding}: \code{nn.Linear(patch\_dim, embed\_dim)} nella variante classica,
\classe{QuantumPatchEmbedding} in quella quantistica. Il percorso quantistico
segue la pipeline della \cref{eq:pipeline-vqc}, con \code{classical\_pre\_process}
a comprimere $147 \to 8$ prima del circuito e $\nq = 8$ qubit. L'encoder è un
\code{nn.TransformerEncoder} standard, identico nelle due varianti: l'attenzione
non è quantizzata (\cref{ssec:precisazione-02}).

L'appaiamento segue la \cref{sec:multiseed}. Una ripresa da checkpoint non
reimposta l'RNG globale: ricarica dal checkpoint gli stati di PyTorch e del
generatore del DataLoader.

Da \file{params.json} risultano \num{3169} parametri totali nella variante
classica e \num{3217} in quella quantistica. La differenza di 48 corrisponde
esattamente ai pesi variazionali, $\nqlayers \cdot \nq \cdot 3 = 2 \cdot 8 \cdot 3$,
perché la compressione $147 \to 8$ è presente e identica in entrambe le varianti,
come \code{nn.Linear} nella prima e come \code{classical\_pre\_process} nella
seconda.

A differenza dell'esperimento~03, il 02 non produce un conteggio per componente:
\file{params.json} contiene solo \code{total} e \code{trainable}, e la
scomposizione riportata sopra è ricostruita per calcolo anziché letta da un file.

\begin{limite}
Le run storiche dell'esperimento~02 non sono utilizzabili per il confronto
principale. Alcune confrontano dataset o configurazioni differenti, mentre la
precedente run quantistica su PathMNIST si interrompe dopo otto delle cinquanta
epoche previste; inoltre il vecchio driver non garantiva né la stessa
inizializzazione dei tensori condivisi né lo stesso ordine dei minibatch per ogni
coppia. Le conclusioni su H2 poggiano esclusivamente sulle run della campagna
definitiva.
\end{limite}

Il test primario è quindi limitato a PathMNIST e cinque seed appaiati, condizioni
più ristrette di quelle di H3. Eventuali repliche su altri dataset restano analisi
secondarie e non cambiano retroattivamente il criterio di decisione primario.

\section{Esperimento 03 --- Ablazione sulla regolarizzazione}
\label{sec:esp-ablazione}

È l'unico dei quattro in cui il componente quantistico resta isolato, senza
confusione con uno strato classico addestrabile a monte, e riceve per questo una
descrizione più estesa.

\begin{ipotesi}[H3]
Il circuito quantistico variazionale agisce da regolarizzatore implicito: a parità
di corpo del modello e di strato di compressione, il gap di generalizzazione della
variante \code{quantum\_reg} è inferiore a quello della variante
\code{bounded\_mlp}.
\end{ipotesi}

Metrica primaria è $\gap = \acctr - \accte$. Il confronto primario è
\code{quantum\_reg} contro \code{bounded\_mlp}, per le ragioni esposte nella
\cref{sec:baseline}; \code{vanilla} resta riferimento secondario.

Tutte e tre le varianti condividono la \emph{stessa} \code{nn.Linear(147, 10)} di
compressione e differiscono unicamente in ciò che avviene dopo di essa.

\begin{table}[htbp]
\centering
\caption[Le tre varianti dell'ablazione]{%
Le tre varianti dell'esperimento~03 e i rispettivi conteggi di parametri, letti
da \file{params.json} (PathMNIST, 9 classi). Encoder \num{2660}, classificatore
$99$ e altri parametri $200$ sono \emph{identici} nelle tre varianti.}
\label{tab:varianti-ablazione}
\small
\begin{tabularx}{\textwidth}{@{}l X r r@{}}
\toprule
\textbf{Variante} & \textbf{Post-compressione} & \textbf{Patch embed} & \textbf{Totale} \\
\midrule
\code{vanilla} & nessuna & \num{1480} & \num{4439} \\
\code{bounded\_mlp} & $\tanh\cdot\pi \to$ MLP($10\to20\to10$, senza bias, GELU) $\to \tanh$ & \num{1880} & \num{4839} \\
\code{quantum\_reg} & $\tanh\cdot\pi \to$ VQC($\nqlayers=2$, $\nq=10$) $\to \expval{Z}$ & \num{1540} & \num{4499} \\
\bottomrule
\end{tabularx}
\end{table}

Le tre varianti costituiscono i tre livelli di una scomposizione, ed è questo a
rendere l'ablazione informativa. Il
passaggio da \code{vanilla} a \code{bounded\_mlp} isola l'effetto del
\textbf{bounding}, perché si aggiunge la limitatezza dell'output e si aggiunge
capacità, ma non unitarietà né entanglement. Il passaggio da \code{bounded\_mlp} a
\code{quantum\_reg} isola invece l'effetto di \textbf{unitarietà ed
entanglement}: la limitatezza dell'output è la stessa, la capacità parametrica del
blocco è addirittura inferiore, e la sola differenza strutturale è la natura
quantistica della trasformazione.

L'ipotesi H3 riguarda il secondo passaggio. Il primo funziona da effetto di
controllo e serve a stabilire che l'esperimento sia sensibile: se nemmeno
l'aggiunta del bounding producesse un effetto misurabile sul gap, l'esperimento
non avrebbe la risoluzione necessaria per rilevare l'effetto cercato, e questo va
constatato prima di interpretare il secondo confronto.

I flag \code{log\_gradients} e \code{log\_activations} sono attivi, e gli
\emph{hook} descritti nella \cref{ssec:hook} sono agganciati a
\code{patch\_embed.compression}, che è lo stesso modulo in tutte e tre le
varianti: stessa classe, stessa forma, stessa posizione nel grafo. È questa
identità a rendere confrontabili le tracce di norma del gradiente e di frazione di
saturazione fra varianti, e non soltanto entro una variante. La campagna primaria
non comprende invece la valutazione di robustezza al rumore, esclusa per non
ampliare a posteriori la famiglia di endpoint.

La campagna comprende $3 \times 5 \times 3 = 45$ run, tutte completate. Ogni run
contiene 400 record di storia e 11 valutazioni di validation (epoca 1 e poi ogni
40 step), e il test è valutato soltanto dopo la selezione del checkpoint sulla
validation. Per ciascuna vengono registrati accuracy di training e di test, gap,
macro-AUC, macro-F1, loss di test, conteggio dei parametri per componente, storia
per epoca, tempi, configurazione e manifest.

Il file \file{main.py} non calcola alcun test statistico: l'aggregazione fra seed e
i confronti appaiati avvengono in un passo separato che invoca
\code{paired\_comparison()} sui \file{results.json} raccolti. La separazione fra
esecuzione e analisi è deliberata, perché consente di rieseguire l'analisi senza
rieseguire l'addestramento.

\section{Esperimento 04 --- Kernel di fedeltà quantistica}
\label{sec:esp-kernel}

\begin{ipotesi}[H4]
Il kernel di fedeltà quantistica separa le classi meglio della similarità coseno
sullo stesso spazio di feature compresso: $F_q > F_c$.
\end{ipotesi}

Metrica primaria è il criterio di Fisher della \cref{eq:fisher}; il rapporto
intra/inter $\rho$ è secondario, con l'avvertenza sulla non confrontabilità delle
scale posta nella \cref{sec:metriche}.

L'esperimento non addestra nulla: confronta due funzioni di similarità su dati
fissati. È un pregio metodologico, perché nessuna delle conclusioni può essere
attribuita a una scelta di iperparametri di ottimizzazione, a un'inizializzazione
sfortunata o a un numero insufficiente di epoche, cioè alle spiegazioni
alternative disponibili per gli altri tre esperimenti.

Da BreastMNIST si estraggono 80 campioni di classe 0 (\emph{malignant}) e 20 di
classe 1 (\emph{normal/benign}), scelti con
\code{np.random.default\_rng(seed).choice(..., replace=False)}. Le immagini sono
appiattite e divise per 255, standardizzate con \code{StandardScaler}, compresse
con \code{PCA(n\_components=8, svd\_solver="full")} e riscalate in $[0,\pi]$ con
\code{MinMaxScaler}. Il solver completo elimina la casualità implicita che
\code{svd\_solver="auto"} introduceva scegliendo la PCA randomizzata; indici
selezionati, solver e varianza spiegata sono salvati in \file{results.json}.

La composizione 80/20 inverte la prevalenza del training set originale (147/399).
Per separare il risultato dalla scelta di campionamento, la stessa analisi è
ripetuta con composizione bilanciata 50/50 e con la prevalenza approssimata 27/73.
La compressione è di $784 \to 8$ componenti, cioè un fattore 98, ed è identica per
i due kernel: entrambi operano sulla stessa informazione residua, punto ripreso
nella \cref{sec:compressione}.

\begin{definizione}[Kernel di fedeltà]
Dato un circuito di codifica $\Uenc$,
\begin{equation}
  K(x,y) \;=\; \abs{\bra{0}\Uenc^{\dagger}(y)\,\Uenc(x)\ket{0}}^2 ,
  \label{eq:fidelity-kernel}
\end{equation}
cioè la probabilità di misurare lo stato $\ket{0\ldots0}$ dopo aver codificato $x$
e decodificato $y$. È simmetrico, vale $1$ sulla diagonale e assume valori in
$[0,1]$.
\end{definizione}

Il circuito definitivo usa \code{qml.IQPEmbedding} con due ripetizioni ed
entangler $ZZ$ fra qubit adiacenti; i termini di fase a due corpi codificano i
prodotti $x_i x_j$ come nella famiglia proposta da
\textcite{havlicek2019supervised}. Il valore del kernel è \code{probs[0]} e nessun
parametro è addestrabile: il kernel è funzione dei soli dati. La baseline
preregistrata è la similarità coseno fra i vettori riscalati in $[0,\pi]$ e
normalizzati a norma unitaria, $S = \hat{X}\hat{X}^{\top}$; la controvalidazione
aggiunge un kernel RBF sulle componenti PCA centrate, con $\gamma=1/(2m)$ e $m$
distanza quadratica mediana fra le coppie di training, scelta che non usa le
etichette.

Una versione precedente del repository applicava Hadamard, $R_Z$, una cascata di
CNOT e un secondo $R_Z$. Quel circuito produce un kernel di fedeltà valido, ma non
è IQP, perché la cascata CNOT non è diagonale e mancano i termini $ZZ$ di prodotto
fra feature. È conservato come \code{custom\_cnot\_rz} per auditabilità e non
fornisce i risultati definitivi.

La campagna usa 20 ripetizioni con seed 42--61; ciascuna impiega il seed
$\code{seed} + r$ e quindi un sottocampione distinto. I sottocampioni possono
sovrapporsi e non costituiscono 20 dataset indipendenti: l'unità di variazione è
la selezione del training set, mentre il test ufficiale della controvalidazione
predittiva rimane fisso. Il calcolo pairwise richiede $n(n+1)/2 = \num{5050}$
valutazioni di circuito per 100 campioni, con crescita quadratica nella
numerosità, e sulla CPU di riferimento una ripetizione richiede in media
\SI{38,24}{\second}. Un controllo indipendente costruisce invece i 100 vettori di
stato e ne calcola il Gram: sul simulatore fornisce la stessa matrice entro
$2{,}0\times10^{-15}$, pur non rappresentando una procedura disponibile su
hardware quantistico.

Il verdetto \code{within\_10\_percent} prodotto dal codice è descrittivo e non è un
test statistico: confronta i due valori di Fisher di una singola ripetizione con
una soglia arbitraria del \SI{10}{\percent}. La conclusione statistica su H4
proviene esclusivamente da \code{paired\_comparison()} applicata alle ripetizioni.
